% Transcription — fonds Alexandre Grothendieck, shelfmark 34, batch 7 (pages 121-140).
% Established from the facsimile archives/batches/34/batch-07.pdf (a mirror of
% https://grothendieck.umontpellier.fr/34.pdf), per the skill
% transcribe-grothendieck.
% Two runs of a manuscript draft for SGA 7. Pages 121-127 are his pagination
% 37-43: the end of section 4 (tame cohomology of the Kummer tower U' over a
% normal-crossings complement, the subgroup G' acting trivially, 4.15-4.17),
% section 5 (proof of the monodromy theorem via the spectral sequence (5.1))
% and section 6 (application to tame vanishing-cycle sheaves, theorem 6.2,
% remarks 6.3-6.4). Page 128 is blank and skipped. Pages 129-140 are his
% pagination V.1-V.13 (131 carries V.3 with a pasted slip V.4): a chapter
% headed IV, "Propriétés générales des faisceaux de cycles évanescents" —
% fibres of the vanishing-cycle sheaves, vanishing bounds from the affine
% Lefschetz theorem (2.1-2.8) and from the global Lefschetz theorem (3.1-3.2).
% Pass: Opus 5.5 (claude-opus-5-5), 2026-10-03 — first pass, unchecked against the pages by a human.
%
\documentclass[11pt,a4paper]{article}
\input{../preamble/grothendieck}

\folder{34}
\batch{7}
\pages{121}{140}
\foldertitle{SGA 7 : notes manuscrites (s.d.).}
\dating{[vers 1967-1973]}
\watermark{Édition de démonstration}

\begin{document}

\page{121}
\note{p. 37 de l'auteur. La page commence au milieu d'un argument venu de la page précédente : le revêtement $U'$ et les numéros (4.8.1), (4.14.5) renvoient au début du § 4, hors de ce lot.}

Ce revêtement est de la forme
\[
U' = U(n) \times^{G} \Gamma ,
\]
pour un homomorphisme $\varphi : G \to \Gamma$ dont le calcul est immédiat, et qui \struck{\ill{}} s'explicite, moyennant (4.8.1) et (4.14.5), comme
\[
\text{(4.14.6)} \qquad \varphi(\lambda_1, \ldots, \lambda_r) = \sum_{1}^{r} m_i \lambda_i .
\]
On en conclut moyennant (4.10.1) et (4.13.1) le calcul de $H^*(U', \Lambda)$. On trouve en particulier, grâce à 4.11 :

\emph{Corollaire} 4.15. \add{Avec les notations de 4.14.} Soit $m$ le pgcd des $m_i$, \struck{\ill{}} \struck{et} $m'$ le morceau de $m$ premier à $p$, et $\Gamma'$ \add{$= m'\Gamma = m\Gamma$} le sous-groupe de $\Gamma \simeq \prod_{\ell \neq p} \mathbf{Z}_\ell(1)$ d'indice $m'$. Alors le sous-groupe des éléments de $\Gamma$ qui opèrent trivialement sur $H^*(U', \Lambda)$ (où $\Lambda = \mathbf{Z}/n\mathbf{Z}$, $n \geqslant 1$, $(n,p) = 1$) est égal à $\Gamma'$. \struck{Pour tout entier premier $\ell$}

On notera que $\Gamma'$ est \emph{\uncertain{indépendant}} de $\Lambda$, i.e. de $n$.

\struck{4.16. Nous allons mettre \ill{} \ill{} limites sur l'\ill{} \ill{}}
\note{deux lignes et un pavé barrés, en grande partie illisibles.}

4.16. Nous allons mettre les résultats 4.14, 4.15 sous une forme \struck{\ill{}} plus globale, qui nous sera utile. Pour ceci, supposons donné un schéma \add{régulier} $X$, un diviseur $D$ sur $X$ à croisements normaux, et une section
\[
t \in \Gamma(X, \mathcal{O}_X) .
\]
\struck{Supposons que $\pi(t)$ \ill{}} \struck{Soit $U =$}

\page{122}
\note{p. 38 de l'auteur.}

Soit
\[
U = X - \operatorname{supp} D ,
\]
et supposons que \struck{\ill{}}
\[
t|U \in \Gamma(U, \mathcal{O}_U^*) , \quad \text{i.e.} \quad \operatorname{supp} \operatorname{div} t \subset \operatorname{supp} D .
\]
\struck{Posons, pour tout entier $n$ premier aux caractéristiques,} Supposons pour simplifier \struck{qu'il \ill{} $X$ \ill{} \ill{}} \struck{\add{$X$ de car. nulle, ou de car. $p$ non nulle,}} \struck{qu'il existe au plus un nombre premier $p$ qui ne soit pas inversible sur $X$.} \uncertain{que}
$V(t)$ soit de car. $p > 0$ et que tous les premiers $\ell \neq p$ soient inversibles sur $X$. Pour tout \struck{entier premier} entier $n > 0$ premier aux car. résiduelles de $X$, \struck{\ill{}} soit
\[
\begin{cases}
X'(n) = \operatorname{Spec} \mathcal{O}_X[T]/(T^n - t) , \\
U'(n) = X'(n) \times_X U ,
\end{cases}
\]
de sorte que les $U'(n)$ forment un système projectif de rev. \struck{étales \ill{}} principaux de $U$, de groupes les $\mu_{n,U}$. On considère
\[
U' = \varprojlim U'(n) ,
\]
qui est principal sur $U$ de groupe \struck{$\mu_n$} $G_U$, où
\[
G = \varprojlim \mu_n ,
\]
et on considère les \struck{projections} \add{morphismes canoniques}

\begin{tikzcd}
U' \arrow[r] \arrow[rr, bend left=30, "f"] & U \arrow[r] & X .
\end{tikzcd}

Soit $n \geqslant 1$ premier aux car. résiduelles de $X$, et \struck{$\Lambda = \mathbf{Z}/n\mathbf{Z}$}
\[
\Lambda = \mathbf{Z}/n\mathbf{Z} ,
\]
on s'intéresse aux \add{faisceaux}
\[
R^i f_*(\Lambda_{U'})|X_0 , \qquad X_0 = V(t)
\]

\page{123}
\note{p. 39 de l'auteur.}

sur $X_0$, et à l'opération de $G|X_0$ sur ces faisceaux. Pour toute composante irréductible $D_i$ de $D$, soit $m_i$ la multiplicité de \struck{\ill{}} $D_i$ dans $\operatorname{div}(\pi)$ ; soit $m$ le pgcd des $m_i$, et soit $m' = m$ si $X$ de car. nulle, $m' =$ partie de $m$ première à $p = \operatorname{car.} V(\pi)$ dans le cas contraire.
\note{ici l'auteur écrit $\pi$ là où la page précédente a $t$.}
Écrivons $G$ additivement et considérons le groupe profini
\struck{$m'G$}
\[
G' = m'G \;(= mG) \subset G
\]
sur $X$. \struck{\ill{}}

\emph{Corollaire} 4.17. Avec les \add{hypothèses et notations} \struck{hypothèses} de 4.16, supposons de plus que pour tout localisé strict $\overline{X}^0$ de $X$ en un point de $X_0$, les $\overline{X}(n)$ définis en termes du diviseur induit $\overline{D}$ (cf. 4.8) satisfont aux th. de pureté. Alors l'opération de $G'|X_0$ sur \struck{\ill{}} $R^*f_*(\Lambda_{U'})|X_0$ est triviale.

Cela résulte aussitôt de 4.15 par localisation stricte.

\struck{4.18. En fait, on peut préciser la structure de} \struck{$R^*f_*(\Lambda_{U'})|X_0$,} \struck{\uncertain{comme}} \struck{\uncertain{faisceau}} \struck{\uncertain{d'anneaux}} \struck{\uncertain{gradués}} \struck{\ill{} \ill{}, dont} \struck{\ill{} \ill{} d'abord} \struck{(4.18.1) $R^*f$}
\note{paragraphe 4.18 encadré et barré de traits obliques.}

\page{124}
\note{p. 40 de l'auteur.}

\section*{5. Démonstration du théorème de monodromie.}

Nous prouvons maintenant 3.2, en utilisant la suite spectrale (2.5 \add{bis de 2.10}) (pour $\nu$ variable). Les résultats des n° \struck{\ill{}} \add{précédents} prouvent que les $(\underline{\Psi}^q_\nu)$\add{$_{\mathrm{tame}}$} sont des faisceaux constructibles, donc les $E_2^{pq}$ \struck{\ill{}} (pour $\nu$ fixé) sont des groupes finis. Donc, passant à la limite sur $\nu$, on trouve une suite spectrale, compatible avec les opérations de $I_t$ et de $\pi$ \ill{} \ill{}
\[
\text{\struck{(5.1)}} \quad \text{\struck{$H^*(\overline{U}, \mathbf{Z}_\ell)^P \Longleftarrow H^p(\widetilde{X}_0, (\underline{\Psi}^q)_{\mathrm{tame}})$}}
\]
\struck{avec}
\[
\text{\struck{(5.2)}} \quad \text{\struck{$(\underline{\Psi}^q)_{\mathrm{tame}} = \varprojlim_\nu (\underline{\Psi}^q_\nu)_{\mathrm{tame}}$}}
\]
\[
\text{(5.1)} \qquad H^*(\overline{U}, \mathbf{Z}_\ell)^P \Longleftarrow E_2^{pq} = \varprojlim_\nu H^p(\widetilde{X}_0, (\underline{\Psi}^q_\nu)_{\mathrm{tame}}) .
\]
On a vu dans 4.17 que le s-groupe $I'_t \subset I_t$ opère trivialement sur les $\underline{\Psi}^q_\nu$ donc \struck{\ill{}} sur les $E_2^{pq}(\nu) = H^p(\widetilde{X}_0, \underline{\Psi}^q_\nu)$, donc sur leur limite $E_2^{pq}$. Donc $I'$ opère trivialement sur tous les \struck{\ill{}} termes de la suite spectrale (5.1), donc sur le \struck{\ill{}} gradué associé à l'aboutissement. Comme le \struck{filt.} gradué, en \uncertain{degré} $i$, \uncertain{a au plus} $i+1$ facteurs, on a gagné.

\page{125}
\note{p. 41 de l'auteur.}

\section*{6. Application aux faisceaux de cycles évanescents.}

6.1. Soit à nouveau $X$ de type fini sur $S$, \add{trait,} strictement local, et soit $U$ un ouvert de $X_\eta$, que nous supposons \emph{régulier}. \add{Soit $Y = X - U$, et} Supposons qu'il existe un morphisme propre $p : X' \to X$, \struck{tel que} avec $X'$ régulier, $p^{-1}(Y)$ un diviseur à croisements normaux, et $U' = p^{-1}(U) \to U$ un isomorphisme.
\note{en marge gauche, un diagramme, transcrit tel qu'il est écrit (les deux sommets de droite portent $U$).}

\begin{tikzcd}
X' \arrow[d, "p"'] & U \arrow[l, hook'] \arrow[d, "\wr"] \\
X \arrow[d, "f"'] & U \arrow[l, hook'] \\
S &
\end{tikzcd}

(On sait par Hironaka [\,] qu'il en est ainsi si $S$ est excellent de car. nulle, et \struck{sans doute} il semble plausible qu'il en est de même dès que $S$ est excellent, sans restriction de caractéristique. D'ailleurs l'hypothèse « $S$ excellent » est anodine, car on s'y ramène pour les questions qui nous intéressent par passage au complété). \struck{Sous ces conditions} Choisissons un entier $\ell \neq p$, et considérons, pour \struck{tout} \add{un} entier $q$ \add{fixé}, les faisceaux
\[
\text{(6.1.1)} \qquad (\underline{\Psi}^q_\nu)_{\mathrm{tame}} = (\underline{\Psi}^q_\nu)^P
\]
(invariants \struck{par} \add{dans $\underline{\Psi}^q_\nu$ sous l'action du} le $p$-groupe de Sylow $P$ de $I$) \struck{\ill{}}, \struck{\ill{}} \add{(cf. 2.6 bis)}, relatifs aux coefficients
\[
\Lambda_\nu = \mathbf{Z}/\ell^{\nu+1}\mathbf{Z} \qquad (\nu \geqslant 0) .
\]
Ils forment un système projectif de faisceaux de torsion sur $X_0$.

\emph{Théorème} 6.2. Avec les notations précédentes :

(i) Les faisceaux $(\underline{\Psi}^q_\nu)_{\mathrm{tame}}$ sont constructibles et \struck{ils} le système projectif qu'ils forment \add{est AR-$\ell$-adique (SGA 5 V 3.2.2)}, \struck{\add{\ill{}} satisfait à la condition MLAR (SGA 5 V 2.1.1)}, donc définissent un faisceau $\ell$-adique constructible (SGA 5 VI 1.1)
\[
(\underline{\Psi}^q)_{\mathrm{tame}} = \varprojlim_\nu (\underline{\Psi}^q_\nu)_{\mathrm{tame}} .
\]

\page{126}
\note{p. 42 de l'auteur. En tête, à droite, une formule barrée : \struck{$(n = q(q+1))$}.}

(ii) Il existe un sous-groupe \add{ouvert} $I'$ de $I$, \struck{et un entier $n$,} tel que pour tout $q$ et $\nu$, l'action de $I'$ sur $(\underline{\Psi}^q_\nu)_{\mathrm{tame}}$ soit unipotente d'échelon \struck{\ill{}} $\leqslant q+1$. \struck{Alors} \add{Alors} l'action de $I'$ sur $(\underline{\Psi}^q)_{\mathrm{tame}}$ est également unipotente d'échelon \struck{\ill{}} $\leqslant q+1$.

\emph{Démonstration.} Soient $g, g'$ les morphismes canoniques de $U_{\overline{K}_t}$ dans $X'$, $X$. On a

\begin{tikzcd}
X' \arrow[d, "p"'] & \\
X & U_{\overline{K}_t} \arrow[l, "g"] \arrow[lu, "g'"']
\end{tikzcd}

\[
\text{(6.2.1)} \qquad (\underline{\Psi}^q_\nu)_{\mathrm{tame}} = R^q g_*(\Lambda_\nu)|X_0
\]
et \struck{\uncertain{permet}} d'autre part la suite spectrale de Leray des morphismes composés $g = p g'$ donne
\[
\text{(6.2.2)} \qquad R^* g_*(\Lambda_\nu) = (\underline{\Psi}^*_\nu)_{\mathrm{tame}} \Longleftarrow E_2^{pq}(\nu) = R^p p_*(R^q g'_*(\Lambda^q_\nu)) .
\]
\note{au-dessus de $E_2^{pq}(\nu)$, un premier essai biffé et illisible.}
En vertu des \struck{\ill{}} \add{calculs du § 4}, il existe un s-groupe ouvert $I'$ de $I$ qui, \add{pour tous $q, \nu$,} opère \add{trivialement sur les} \struck{sur les faisceaux unipotents d'échelon} $R^q g'_*(\underline{\Psi}^q_\nu)$. Il s'ensuit que $I'$ opère \struck{\ill{}} \add{trivialement} \struck{\uncertain{unipotentes} $q+(0+1)+(1+1)+\cdots+($} sur \struck{les faisceaux unipotents d'échelon} les $E_2^{pq}(\nu)$, donc \struck{sur} il opère de façon unipotente \add{d'échelon $q+1$} \struck{d'échelon} sur $(\underline{\Psi}^q_\nu)_{\mathrm{tame}}$ \struck{d'échelon $(0+1)+(1+1)+\cdots+(q+1)$} \struck{$= q(q+1)/2$}, ce qui prouve la première assertion de (ii).
\note{un calcul encadré et barré : la somme $(0+1)+(1+1)+\cdots+(q+1) = q(q+1)/2$ ; le cadre renvoie à « d'échelon $q+1$ ».}
D'autre part, les calculs explicites du § 4 montrent que les $R^q g'_*(\Lambda^q_\nu)|X_0$ sont des faisceaux \emph{constructibles}, qui pour $\nu$ variable forment un faisceau \struck{\ill{}} $\ell$-adique (i.e. $R^q g'_*(\Lambda_{\nu'})|X'_0 \simeq R^q g'_*(\Lambda_\nu)|X'_0 \otimes_{\Lambda_\nu} \Lambda_{\nu'}$ si $\nu \geqslant \nu'$). En vertu de SGA 5 VI 2.2, il s'ensuit que

\page{127}
\note{p. 43 de l'auteur.}

les $R^p p_*(R^q g'_*(\Lambda_\nu))|X_0$ \add{$= R^p p_*(R^q g'_*(\Lambda_\nu)|X'_0)$ (th. de changement de base propre)}, pour $\nu$ variable, forment un système AR-$\ell$-adique de faisceaux sur $X$, donc il en est de même des $R^i g_*(\Lambda_\nu)|X_0 = (\underline{\Psi}^i_\nu)_{\mathrm{tame}}$, en vertu de (SGA 5 V 5.2.1 et 5.2.4), d'où (i).

La deuxième assertion de (ii) \add{relative à $(\underline{\Psi}^q)_{\mathrm{tame}}$ lui-même} est une conséquence immédiate de la première, et de ce qui précède.

\emph{Remarques} 6.3. Si $m'$ est l'entier défini dans 3.2 en termes des multiplicités des composantes de $X'_0$ dans $X'$, alors la démonstration précédente prouve que dans 6.2 on peut prendre pour $I'$ le s-groupe \add{ouvert} d'indice $m'$ de $I$.

6.4. Il est plausible qu'on devrait plus raisonnablement pouvoir démontrer des résultats analogues à 6.2, mais sans supposer $U$ \emph{régulier}, et pour les faisceaux $\underline{\Psi}^q_\nu$ eux-mêmes, au lieu des sous-faisceaux $(\underline{\Psi}^q_\nu)_{\mathrm{tame}}$. Pour ce dernier point, il semble qu'il faille cependant \add{d'abord} \struck{\ill{}} \add{résoudre par l'affirmative} \struck{\ill{}} la question suivante : \struck{\ill{}} Supposons \emph{remplies les conditions de 3.2}, \emph{\uncertain{ici les composantes} \struck{\ill{}} \uncertain{réduites} de $X_0$ lisses sur $\overline{k}$}, \emph{a-t-on} $(\underline{\Psi}^q_\nu)_{\mathrm{tame}} = \underline{\Psi}^q_\nu$, i.e. \emph{$P$ opère-t-il trivialement sur les $\underline{\Psi}^q_\nu$ ?} \struck{\ill{}}
\marginal{et les multiplicités de $X_0$ premières à $p$}

Cela impliquerait bien sûr, \add{si $X$ est propre sur $S$,} que $P$ opère trivialement sur les $H^i(U_{\overline{\eta}}, \Lambda_\nu)$, chose qui n'est pas non plus prouvée à l'heure actuelle. \struck{\ill{}} Pour le cas $i = 1$, c'est vrai, et ce sera démontré dans un exposé ultérieur.
\struck{Mais prenons, \ill{} $\dim X_0 = \dim X_\eta = 1$, $U = X_\eta$, $q = 1$, la réponse \ill{} \ill{} \ill{} \ill{} \ill{} pas \ill{}}
\note{l'argument se poursuit au-delà de cette page, hors du lot ; la page 128 est blanche.}

\page{129}
\note{p. V.1 de l'auteur. Une nouvelle rédaction commence ici, sous le numéro IV ; le titre a été corrigé en surcharge, d'une autre encre.}

\section*{IV \struck{\ill{} Cas des points} \add{Propriétés générales des faisceaux de cycles évanescents.} \struck{singuliers isolés de la fibre spéciale.}}

1. Calcul \struck{des} \add{de \struck{\ill{}} la} fibre des faisceaux des cycles évanescents, \add{comme} \struck{\uncertain{terme} de la} cohomologie de la fibre générique géométrique d'un localisé strict.

\marginal{N.B. \uncertain{Omettre toujours} les $\nu$ en indice aux $\Lambda_\nu$, sauf en cas de confusion possible ; \uncertain{pourra de même} laisser tomber $\Lambda$ dans $H^*(X, \Lambda)$ etc., et écrire $H^*(X)$ …}

\page{130}
\note{p. V.2 de l'auteur. Au-dessus de $X'$, $f'$, $U'$, un signe (tilde ?) a été biffé à chaque occurrence ; on transcrit sans lui.}
\marginal{dans Exp. V}

\struck{1.1} 1.1. Nous allons donner une interprétation frappante des fibres $(\underline{\Psi}^{(q)}_\nu)_{\overline{x}}$ des faisceaux introduits dans (\add{III} 2.6), où $\overline{x}$ est un point géométrique de $\widetilde{X}_0$. Pour ceci, introduisons \add{nous pouvons supposer} le localisé strict $X'$ de $\widetilde{X}$ en $\overline{x}$ ; \struck{\ill{}} le localisé strict correspondant de $S$ égal à $\widetilde{S}$, de sorte qu'on a un morphisme canonique
\[
\text{\struck{(1.1.1)}} \quad \text{(1.1.1)} \qquad f' : X' \longrightarrow \widetilde{S}
\]
\marginal{faire $\widetilde{S} = S$ ici !}
Si $U'$ est l'image inverse de $U$ dans $X'$, c'est un ouvert de la fibre générique $X'_{\widetilde{\eta}}$ de $f'$ sur $\widetilde{S}$ :
\[
\text{(1.1.2)} \qquad U' \subset X'_{\widetilde{\eta}} .
\]
Ceci posé, le \struck{\uncertain{calcul des}} calcul habituel des fibres de faisceaux \add{$R^i g_*$} (SGA 4 VII) montre que

\marginal{laisser tomber $U$ !}

\page{131}
\note{p. V.3 de l'auteur ; dans le bas du feuillet est collé un papillon numéroté V.4 par l'auteur (131A pour les archivistes), transcrit à la suite. Dans le diagramme, les signes au-dessus de $X'$, $X'_{\overline{\eta}}$, $U'$ sont biffés, comme à la page précédente.}

\begin{tikzcd}[column sep=small, row sep=small, nodes={font=\scriptsize}]
 & & X' \arrow[r, hook'] \arrow[dl] & X'_{\overline{\eta}} \arrow[dl] & U' \arrow[l, hook'] \arrow[dl] & & & \overline{U}' \arrow[lll] \arrow[dl] \\
X \arrow[d] & \widetilde{X} \arrow[l] \arrow[d] & \widetilde{X}_{\widetilde{\eta}} \arrow[l, hook'] \arrow[d] & \widetilde{U} \arrow[l, hook'] & & & \overline{U} \arrow[lll] \arrow[d] \arrow[llll, bend left=20, "g"'] & \\
S & \widetilde{S} \arrow[l] & \widetilde{\eta} \arrow[l] & & & & \overline{\eta} \arrow[llll] &
\end{tikzcd}

\note{diagramme redessiné de l'esquisse : les positions sont approximatives ; la flèche verticale sous $\widetilde{U}$ et la flèche $\overline{\eta} \to \widetilde{\eta}$ (corrigée en surcharge) sont lues avec doute.}

un isomorphisme canonique
\[
(\underline{\Psi}^{(q)}_\nu)_{\overline{x}} \simeq \text{\struck{$H^q(\overline{U}' \times_{\widetilde{X}} \overline{U})$}} \; H^q(\overline{U}', \Lambda_\nu) , \quad \text{avec} \quad \overline{U}' \simeq X' \times_{\widetilde{X}} \overline{U} ,
\]
et comme on a
\[
X' \times_{\widetilde{X}} \overline{U} \simeq U' \times_{\widetilde{U}} \overline{U} \simeq U' \times_{\widetilde{\eta}} \overline{\eta} ,
\]
on trouve finalement
\[
\text{(1.1.3)} \qquad (\underline{\Psi}^{(q)}_\nu)_{\overline{x}} \simeq H^q(\overline{U}', \Lambda_\nu) , \qquad \overline{U}' = U' \otimes_{\widetilde{\eta}} \overline{\eta} = U' \otimes_{\widetilde{K}} \overline{K} .
\]
Si par exemple $U = X_\eta$ (cas de loin le plus important), on trouve
\[
\text{(1.2)} \qquad (\underline{\Psi}^{(q)}_\nu)_{\overline{x}} = H^q(X' \times_{\widetilde{S}} \overline{\eta}, \Lambda_\nu) ,
\]
i.e. \emph{la fibre du faisceau $(\underline{\Psi}^{(q)}_\nu)_{\overline{x}}$ en le point géométrique $\overline{x}$ de $\widetilde{X}$ \struck{\ill{}} n'est autre que le groupe de cohomologie \add{(à coeff. dans $\Lambda_\nu$)} de la fibre géométrique générique sur $\widetilde{S}$ du localisé strict $X'$ de $X$ en $\overline{x}$.} \struck{\ill{} \ill{} \ill{} prouver les relations} \struck{\ill{}} relations (2.8), on \struck{\ill{}} \ill{} compte tenu de \struck{\ill{}} en conclut des expressions correspondantes pour les \add{\emph{fibres génériques des morphismes $X' \to \widetilde{S}$ entre localisés stricts de $X$ et $S$.}}
\marginal{De façon imagée, \uncertain{on peut dire} que les relations globales entre les deux \uncertain{fibres} géométriques $X_{\overline{s}}$, $X_{\overline{\eta}}$ se \ill{} \ill{} des \ill{} \ill{} \ill{} des cycles évanescents. \ill{}}
\note{note marginale écrite en oblique dans le coin inférieur gauche, en partie barrée et en grande partie illisible.}

\note{papillon collé, p. V.4 de l'auteur.}
\struck{2} \struck{Lorsque} Le même argument donne les relations
\[
\text{\struck{(1.1.3)}} \qquad ((\underline{\Psi}^{(q)}_\nu)_{\mathrm{tame}})_{\overline{x}} \simeq H^q(U' \otimes_{\widetilde{K}} \overline{K}_t) ,
\]
en particulier si $U = X_\eta$, cela \struck{devient} prend la forme
\[
\text{(1.3)} \qquad ((\underline{\Psi}^q_\nu)_{\mathrm{tame}})_{\overline{x}} = H^q(X \otimes_{\widetilde{S}} \operatorname{Spec} \overline{K}_t) \; ;
\]
\struck{\ill{}}

\page{132}
\note{p. V.5 de l'auteur.}

\begin{tikzcd}[column sep=small, row sep=small, nodes={font=\scriptsize}]
 & X' \arrow[r, hook'] \arrow[dl] & X'_\eta \arrow[l] \arrow[dl] & U' \arrow[l, hook'] \arrow[dl] \arrow[r, no head] & U'_{\overline{\eta}} \arrow[dl] \\
X \arrow[d, no head] & X_\eta \arrow[l, hook'] \arrow[d] & U \arrow[l] \arrow[d] \arrow[r, no head] & U_{\overline{\eta}} \arrow[d, no head] & \\
S & \eta \arrow[l] & \eta \arrow[l] \arrow[r, no head] & \overline{\eta} &
\end{tikzcd}

\note{diagramme redessiné : la flèche $X_\eta \to U$ est dessinée dans ce sens sur la page, et une mention « (pas ent.) » est inscrite à côté de la flèche verticale sous $X_\eta$ ; plusieurs traits n'ont pas de pointe.}

1.4. Notons que l'on a
\[
X' = \varprojlim_\alpha X_\alpha ,
\]
où les $X_\alpha$ sont étales sur $X$, et affines. On en conclut
\[
U' = \varprojlim_\alpha X_\alpha \times_X U , \qquad U'_{\overline{\eta}} = \varprojlim_\alpha (X_\alpha \times_X U)_{\overline{\eta}} ,
\]
d'où
\[
\text{(1.4.1)} \qquad (\underline{\Psi}^{(i)}_\nu)_{\overline{x}} \simeq H^i(U'_{\overline{\eta}}, \Lambda) \simeq \varinjlim_\alpha H^i((X_\alpha \times_X U)_{\overline{\eta}}, \Lambda) ,
\]
où les \struck{\ill{}} $(X_\alpha \times_X U)_{\overline{\eta}}$ sont des schémas algébriques quasi-affines sur $\overline{\eta}$, déduits par changement de base des schémas quasi-affines $X_\alpha \times_X U \subset X_\alpha$ sur $\eta$.

Lorsque $U = X_\eta$, cela devient

\struck{2. Application}
\[
\text{(1.4.2)} \qquad (\underline{\Psi}^{(i)}_\nu)_{\overline{x}} = H^i(X'_{\overline{\eta}}, \Lambda) \simeq \varinjlim_\alpha H^i(X_{\alpha\overline{\eta}}, \Lambda) ,
\]
où les $X_{\alpha\overline{\eta}}$ sont des schémas \add{algébriques} \emph{affines} sur $\overline{\eta}$ [déduits des schémas affines $X_{\alpha\eta}$ sur $\eta$]. \struck{\ill{}}

N.B. Plus généralement, dans 1.4.1, les $(X_\alpha \times_X U)_{\overline{\eta}}$ sont affines si on suppose que $U \to X$ est affine, i.e. $U \to X_\eta$ est affine.

\page{133}
\note{p. V.6 de l'auteur. Page très reprise : plusieurs ajouts interlinéaires et un passage encadré de traits obliques, qui semble annulé.}

\section*{2. Application du th. de Lefschetz affine.}

2.1. Soit $x \in X_0$, et
\[
\text{(2.1.0)} \qquad n_x = \text{\struck{$\dim X_\eta$}} \; \sup_i \dim X_{i\eta} ,
\]
\add{$X_i$ composantes irréd. de $X$ passant par $x$}. \marginal{Supposons $U \to X_\eta$ affine.} Alors dans (1.4.2) les $X_{\alpha\overline{\eta}}$ sont les affines algébriques \add{sur $\overline{\eta}$} \struck{\ill{} $U_{X_\eta}$} et sont \add{\uncertain{dont les composantes sont au plus de dimension les} $\dim X_{i\eta}$} \uncertain{étales} sur $X_\eta$, donc affines de dim. $\leqslant n_x$, donc \add{(SGA \struck{4} XIV …)} leur $H^i$ est nul pour $i > n_x$. On trouve donc
\[
\text{(2.1.1)} \qquad (\underline{\Psi}^i_\nu)_{\overline{x}} = 0 \quad \text{pour } i > n_x .
\]
On en améliorera ce résultat en prouvant que
\[
\text{(2.1.2)} \qquad (\underline{\Psi}^i_\nu)_{\overline{x}} = 0 \quad \text{pour } i > n_x - \dim \overline{\{x\}} ,
\]
ce qui implique aussi
\[
\text{(2.1.3)} \qquad (\underline{\Psi}^i_\nu)_{\overline{x}} = 0 \quad \text{pour } i > \dim \mathcal{O}_{X_0, x} = \dim_x(X_0) - \dim \overline{\{x\}}
\]
\struck{(\ill{}, dès que \ill{})}.

2.2. Pour ceci, notons que \struck{\ill{}} si $X$ \add{est} plat sur \struck{$X_\eta$} $S$, l'adhérence schématique $\overline{(X_\eta)}_{\mathrm{red}}$ [\uncertain{ainsi du reste que} \struck{\ill{}} les $\underline{\Psi}^{(i)}_\nu$ \add{\uncertain{\ill{}}} \uncertain{sinon} \ill{}], \ill{} \add{dim.} $X_\eta$ \ill{} \ill{} \ill{} $X_0$, \struck{\ill{}} \ill{} \ill{} \struck{\ill{} $n_x = \dim_x X_0$} $n_x = \dim_x X_0$ \struck{\ill{} (2.1.1) et (2.1.2) sont identiques. \ill{} $x$ de $X$, \ill{} \ill{} \ill{}}

\struck{Alors, il existe un morphisme fini $X \xrightarrow{h} Y = S[T_1, \ldots, T_n]$, et si on pose $F = h_*(\Lambda_X)$, $y = h(x)$, on trouve que les $\underline{\Psi}^{(i)}_*(F)$ relatifs \ill{} les faisceaux \ill{} \ill{} \ill{} induits, comme}
\note{ce passage est encadré et barré de deux grands traits obliques.}
$F$ (qui se définissent \ill{} \ill{} \ill{} \ill{} $X$) \ill{}

On \uncertain{peut} supposer \uncertain{aussi} (\uncertain{quitte à} \uncertain{restreindre} $X$) que \struck{$n_x = \dim X_\eta$, Posons $n_x = n$, \uncertain{\ill{}}, soit $n$.} Soit
\[
n = \dim X_\eta = n_x \; ;
\]
\struck{$r = \dim$} soit d'autre part $r = \deg \operatorname{tr.} k(x)/k$.

\page{134}
\note{p. V.7 de l'auteur.}

\struck{et soit $x$}
On peut supposer $X$ affine, et on peut trouver un \uncertain{système} d'éléments
\[
t_1, \ldots, t_r \in \Gamma(X, \mathcal{O}_X)
\]
telles que leurs images dans $k(x)$ forment une base de transc. sur $k$. Les $t_i$ définissent un morphisme
\[
h : X \longrightarrow Y = S[T_1, \ldots, T_r] ,
\]
et si $y = h(x)$, on sait que $k(x)/k(y)$ est fini, donc que $x$ est fermé dans $h^{-1}(y)$. On a donc
\[
\dim_x h^{-1}(y) = \dim \mathcal{O}_{X_0, x} = n - r ,
\]
et quitte à restreindre $X$, on peut donc supposer que les fibres de $X$ sur $Y$ sont de dim. $\leqslant n - r$. Soit $S_1$ le \struck{\ill{}} localisé strict de $Y$ en le point géométrique de $Y$ défini par le pt géom. $\overline{x}$ de $X$, $X_1 = X \times_Y S_1$, $U_1 = U \times_Y S_1$. Alors $X_1$ est affine
\note{un premier diagramme, à gauche, est barré de grands traits obliques ; on transcrit celui qui le remplace, à droite.}

\begin{tikzcd}[column sep=small, row sep=small, nodes={font=\scriptsize}]
 & X' \arrow[r, hook'] \arrow[d] & U' \arrow[d, no head] & U'_{\overline{\eta}} \arrow[l] \arrow[r, no head] \arrow[dd, "\varphi"] & U' \times_{S_1} \overline{\lambda} \arrow[dd] \\
X & X_1 \arrow[l] \arrow[d] & U_1 \arrow[l, hook'] \arrow[d] & & \\
Y = S[T_1, \ldots, T_r] \arrow[d] & S_1 \arrow[l] \arrow[r, no head, "="] & S_1 \arrow[d] & S_{1\overline{\eta}} \arrow[l] \arrow[r, no head] \arrow[d, no head] & \overline{\lambda} \\
S & & S & \overline{\eta} \arrow[l] &
\end{tikzcd}

\note{diagramme redessiné de l'esquisse ; la flèche $Y \to S$ est doublée sur la page, une flèche verticale sous $S_1$ (à gauche du signe $=$) est raturée, et le sommet $\overline{\lambda}$ remplace un symbole biffé ; lectures de $\overline{\lambda}$ douteuses.}

\uncertain{d.f.} sur le trait strict. hensélien $S_1$, de dim. relative $n - r$, $U_1$ en est un ouvert affine, $X'$ est le localisé strict de $X_1$ en le pt géom. $\overline{x}$, $U'$ est l'image inverse de $U_1$ dans $X'$. Soit $L$ le corps des fractions de $S_1$, qui est un

\page{135}
\note{p. V.8 de l'auteur.}

extension \add{algébrique} séparable de $K(T_1, \ldots, T_r)$, donc une extension séparable de $K = k(\eta)$, (et $\lambda = \operatorname{Spec} L \simeq (S_1)_\eta$ \struck{\ill{}}). D'après 2.1.1, appliqué à la situation $X_1 \to S_1$, et $U_1 \subset X_1$, $X' \to X_1$, on sait que
\[
\text{(2.2.1)} \qquad H^i(U' \times_{S_1} \overline{\lambda}, \Lambda) = 0 \quad \text{si } i > n - r , \quad \text{i.e.} \quad R^i \varphi_*(\Lambda_{U'_{\overline{\eta}}}) = 0 \text{ si } i > n - r .
\]
D'autre part, \struck{$\overline{\lambda}$ \ill{}} on sait que \struck{\ill{}} $\operatorname{cd}_\ell$ \uncertain{de} $S_{1\eta} = \lambda$ est \uncertain{$\leqslant r$} pour tous les \uncertain{nombres} premiers $\neq p$ (SGA 4 XV). Dans la suite spectrale
\marginal{rapprochement, \uncertain{ou} \uncertain{dimension} \uncertain{cohomologique} \uncertain{de} $\operatorname{Gal}(\overline{\lambda}/S_{1\eta})$ \uncertain{qui} \ill{} \ill{} \ill{} …}
\[
H^i(U'_{\overline{\eta}}, \Lambda) \Longleftarrow E_2^{p,q} = H^p(\lambda, R^q \varphi_*(\Lambda_{U'_{\overline{\eta}}}))
\]
\uncertain{dégénère} et donne
\[
H^i(U'_{\overline{\eta}}, \Lambda) \simeq H^0(\lambda, R^i \varphi_*(\Lambda_{U'_{\overline{\eta}}})) \subset H^i(U' \times_{S_1} \overline{\lambda}, \Lambda)
\]
\note{le signe qui précède le dernier terme est lu comme une inclusion, avec doute.}
ce qui implique \uncertain{par} (2.2.1)
\[
H^i(U'_{\overline{\eta}}, \Lambda) = 0 \quad \text{pour } i > n - r ,
\]
et achève la démonstration. On a prouvé :

\emph{Théorème} 2.3. Avec les notations de 2.1, on a \add{pour $x \in X_0$} :
\[
(\underline{\Psi}^q_\nu)_{\overline{x}} = 0 \quad \text{pour } q > n_x - \dim \overline{\{x\}}
\]
\uncertain{(en particulier} \uncertain{pour} $i > \dim \mathcal{O}_{X_0, x} = \dim_x X_x - \dim \overline{\{x\}}$).

De façon équivalente, on a
\[
(\underline{\Phi}^q_\nu)_{\overline{x}} = 0 \quad \text{si } q > n_x + 1 - \dim \overline{\{x\}} , \quad q \geqslant 1 .
\]
\struck{Appliquant le \ill{} \ill{} \ill{} cohomologique} \add{si $n = \dim X_\eta$, $q \geqslant 1$, $X = \overline{X_\eta}$ donc $\underline{\Phi}^0_\nu = 0$ :}
…
\note{passage encadré ; le renvoi « … » semble indiquer que cet encadré s'insère plus haut.}
\add{SGA 4 \uncertain{X}}, on trouve
\[
H^p(X_0, (\underline{\Phi}^q_\nu)|X_0) = 0 \quad \text{si } p \geqslant 2(n + 1 - q) ,
\]
\[
\text{i.e.} \quad p + q \geqslant 2n + 2 - q .
\]
\struck{\emph{Corollaire} 2.3. On a, si $X$ propre sur $S$}

\emph{Corollaire} 2.4. On a
\[
\underline{\Psi}^q_\nu = 0 \quad \text{si } q > \dim X_\eta , \qquad \underline{\Phi}^q_\nu = 0 \quad \text{si } q > \dim X_\eta + 1 .
\]

\page{136}
\note{p. V.9 de l'auteur.}

\struck{N.B.} \struck{Supposons} \add{Considérons} maintenant que l'on \struck{\ill{}} \add{(fermé)} $Z$ \add{de $X_0$} \add{dans} lequel \add{\uncertain{soit}} \add{(III 1. ), dans} $f : X \to S$ est non lisse, \uncertain{\ill{}} \add{d sa} \add{dim.} \struck{\ill{}} \add{\ill{} $\underline{\Phi}^q_\nu \subset Z$}

On a \struck{\ill{}} \add{par SGA 4 X} \ill{}
\[
(*) \qquad H^p(X_0, \underline{\Phi}^q_\nu) \neq 0 \;\; \text{\struck{\ill{}}} \;\; \text{implique} \;\; p \leqslant 2d .
\]
D'ailleurs, on a en vertu de 2.2. (si $n = \dim X_{\overline{\eta}}$)
\[
\delta(\underline{\Phi}^q_\nu) \leqslant n + 1 - q \quad \text{si } q \geqslant 1
\]
donc on trouve \add{$\delta(\underline{\Phi}^q_\nu)$} $\leqslant d$ si $n + 1 - q < d$, i.e. \struck{\ill{}} pour \uncertain{des} tels $q$, \struck{\ill{}} on a par SGA 4 X
\[
q > n + 1 - d \; ; \quad \text{\struck{$2\delta(\underline{\Phi}^q_\nu) =$}}
\]
\[
(**) \qquad H^p(X_0, \underline{\Phi}^q_\nu) \neq 0 \;\; \text{implique} \;\; p \leqslant 2(n + 1 - q) ,
\]
\[
\text{i.e.} \quad p + q \leqslant 2n + 2 - q ,
\]
et \uncertain{on a} \add{pour ces $q$} $2n + 2 - q < 2n + 2 - (n + 1 - d) = n + 1 + d$. D'autre part, si $q \leqslant n + 1 - d$ \struck{et} \add{$H^p(X_0, \underline{\Phi}^q_\nu) \neq 0$ i.e.} $p \leqslant 2d$, on a \uncertain{là encore} $p + q \leqslant n + 1 + d$. Donc on \uncertain{a dans tous les cas} (si $q \geqslant 1$)
\[
H^p(X_0, \underline{\Phi}^q_\nu) \neq 0 \Longrightarrow p + q \leqslant n + d + 1 .
\]
Cette implication est aussi vraie pour $q = 0$, si \uncertain{on} \add{$(d = \ldots \leqslant n = \dim X_0)$} \struck{\ill{}} suppose que $d \leqslant n + 1$ (ce qui est le cas p.ex. si $X = \overline{X_\eta}$, \struck{\ill{}} \uncertain{\ill{}} p.ex. $X$ plat sur $S$]. On trouve \uncertain{\ill{}}
\[
\underline{\Phi}_\nu = 0 \quad \text{pour } i > n + d + 1 .
\]
\note{l'indice manque à $\underline{\Phi}_\nu$ sur la page (ou est surchargé), lu tel quel.}

\emph{Corollaire} 2.5. \add{(2.4.1)} \struck{Supposons} \add{Donc :} \struck{Soit} \add{Soit} $d$ la dimension \struck{\ill{}} \add{\ill{}} \add{\uncertain{du fermé}} de $X_0$ en lequel $f : X \to S$ est non lisse, \add{$n = \dim X_{\overline{\eta}}$, et} \struck{\ill{}} \add{supposons} que $d \leqslant n + 1$ i.e. que $\dim X_0 \leqslant n + 1$ [ce qui est le cas si $X = \overline{X_\eta}$, \uncertain{puisque} $\dim X_\eta = n$]. Alors les groupes $\underline{\Phi}^i_\nu$ des cycles évanescents \add{(III 1. )} à coefficients dans $\Lambda_\nu = \mathbf{Z}/\ell^\nu \mathbf{Z}$, \uncertain{sont} nuls pour \uncertain{les} $i > n + d + 1$, en d'autres termes l'hom.
\marginal{(cel. à $U = X_\eta$)}

\page{137}
\note{p. V.10 de l'auteur.}

\[
H^i(X_0, \Lambda_\nu) \longrightarrow H^i(X_{\overline{\eta}}, \Lambda_\nu)
\]
est un isomorphisme si $i > n + d + 1$, un épimorphisme si $i = n + d + 1$. L'action du groupe d'inertie \struck{\ill{}} $I$ sur $H^i(X_{\overline{\eta}}, \Lambda_\nu)$ est triviale si $i \geqslant n + d + 1$.

Procédant par dualité, on trouve :

\emph{Corollaire} 2.6. Avec les hypothèses précédentes, l'action de $I$ sur $H^i(X_{\overline{\eta}}, \Lambda_\nu)$ est triviale si $i \notin [n - d, n + d]$.

En particulier, si $d = 0$ i.e. si l'ens. $Z$ des pts de $X_0$ où $f$ n'est pas lisse est fini, alors l'action de $I$ sur $H^i(X_{\overline{\eta}}, \Lambda_\nu)$ est triviale si $i \neq n$.

\marginal{interversion avec 2.8}
\emph{Remarque} 2.7. Lorsque $f : X \to S$ est projectif, on peut trouver une démonstration plus simple de \struck{\ill{}} \add{de} 2.5 \uncertain{(ainsi} \add{de} 2.6) \struck{\ill{} \ill{} triviale \ill{}} \struck{\ill{} \ill{} \ill{} \ill{}}. En effet, plongeons $X \subset \mathbf{P}^r_S$, et soit $Y$ une section plane \add{générique} de $X$, obtenue en coupant $X$ par $r - d - 1$ \struck{\ill{}} hyperplans de $\mathbf{P}^r_S$. Si ces derniers sont bien choisis, on peut supposer $Y_0$ \uncertain{lisse} de codim. $d + 1$ dans $X_0$, et alors $Y$ est lisse sur $S$, de codim. $d + 1$ dans $X$. De plus, le th. de Lefschetz affine SGA 4 XIV implique \uncertain{aisément}, \struck{\ill{}} de façon connue, que l'homom. de Gysin
\[
H^i(Y_{\overline{\eta}}, \Lambda_\nu) \longrightarrow H^{i + 2(d+1)}(X_{\overline{\eta}}, \Lambda_\nu(d + 1))
\]

\page{138}
\note{p. V.11 de l'auteur.}

est surjectif si $i \geqslant \dim Y_{\overline{\eta}} = n - d - 1$, \add{et bijectif si $i > n - d - 1$ i.e. si $i + 2(d+1) > n + d + 1$} donc si $i + 2(d + 1) \geqslant n + d + 1$, \struck{Donc le fait que $I$ opère \ill{} sur $H^i(X_{\overline{\eta}}, \Lambda_\nu)$ \ill{} provient du fait qu'il opère trivialement sur $H^i(Y_{\overline{\eta}}, \Lambda_\nu)$, \ill{} est une conséquence bien connue du fait que $Y$ est propre et lisse sur $S$ (SGA 4 XVI \ldots)}
\[
\begin{array}{ccc}
H^i(Y_s, \Lambda) & \longrightarrow & H^{i + 2(d+1)}(X_s) \\
\downarrow{\scriptstyle \varphi_Y} & & \downarrow{\scriptstyle \varphi_X} \\
H^i(Y_{\overline{\eta}}, \Lambda) & \longrightarrow & H^{i + 2(d+1)}(X_{\overline{\eta}})
\end{array}
\]
\note{passage encadré et barré de traits obliques ; à la suite, d'une autre encre : « le diagramme commutatif », puis le carré ci-dessus.}
\marginal{donne alors le résultat voulu ($\varphi_Y$ est bijectif par SGA 4 XVI \ldots)}

\emph{Remarque} 2.8. L'assertion 2.3 suggère la \struck{\ill{}} conjecture suivante. Soit $X = \operatorname{Spec} A$ un schéma strict. \add{local}, \add{de dim. $n$}, $t \in A$, \struck{\ill{}} et $U = X - V(t) = \operatorname{Spec} A_t$, et soit pour tout entier $m$ premier à la car.
\[
X_{(m)} = \operatorname{Spec} A[T]/(T^m - t) , \qquad U_{(m)} = X_{(m)} \times_X U ,
\]
\[
U(\infty) = \varprojlim U(m) .
\]
Est-il vrai que la dim. cohomologique de $U(\infty)$, pour des coefficients premiers à \struck{\ill{}} $p$, \struck{\ill{}} est $\leqslant n - 1$ ? [La conjecture standard SGA 4 XIV \ill{} \struck{\ill{}} \struck{\ill{}} pour $\ell \neq p$ : $\operatorname{cd}_\ell(U(m)) \leqslant n$ pour tout $m$, donc $\operatorname{cd}_\ell(U(m)) \leqslant n$, et non $n - 1$ !]
\note{la fin de la phrase est lue telle qu'elle semble écrite ; la seconde inégalité devrait sans doute porter sur $U(\infty)$.}

\page{139}
\note{p. V.12 de l'auteur.}

\section*{3. Application du th. de Lefschetz \add{(de la section hyperplane)} \struck{\ill{}} global.}

\struck{3.1.} 3.1. Supposons que la profondeur étale \uncertain{cohomologique} \add{pour $\ell$} de $X_s$ et $X_{\overline{\eta}}$ soit $\geqslant m$, i.e. que pour tout $x \in X_s$, on ait
\[
\text{(3.1.1)} \qquad \underbrace{\operatorname{prof\,ét}_x X_s}_{\operatorname{prof\,ét}^*_x(X_s)} + \dim \overline{\{x\}} \geqslant m
\]
et de même pour tout $x \in X_{\overline{\eta}}$,
\[
\text{(3.1.2)} \qquad \underbrace{\operatorname{prof\,ét}_x (X_{\overline{\eta}})}_{\operatorname{prof\,ét}^*_x(X_{\overline{\eta}})} + \dim \overline{\{x\}} \geqslant m ,
\]
où $\overline{x}$ \uncertain{désigne} un \uncertain{point} \uncertain{géométrique} \uncertain{au-dessus de} $x$, \struck{\ill{}} (\uncertain{ainsi} pour \uncertain{alors} …), \uncertain{\ill{}}. On \uncertain{sait que} \ill{} \ill{}
\marginal{cf. SGA 2 XIV … \uncertain{pour la notion de profondeur étale}.}
par l'hypothèse faite sur $X_{\overline{\eta}}$ ou \struck{\ill{}} \add{\ill{}} sur la fibre géométrique $X_{\overline{\eta}}$). \struck{Alors, \uncertain{soit $d$}}

\uncertain{Soit} \add{$d$} comme dans 2.5, et supposons \add{d'abord} $X$ projectif et \emph{plat} sur $S$. Choisissons $Y$ comme dans 2.7, alors le th. de Lefschetz \add{global} \struck{\ill{}} sous la forme SGA 2 XIV 4.6 \emph{entraîne} [\emph{en vertu du changement de base} \struck{\ill{}} \emph{pour les clôtures alg. de $k$ et $K$}] \emph{les relations suivantes} : \uncertain{les}
\[
H^i(X_s, \Lambda) \longrightarrow H^i(Y_s, \Lambda) \qquad (\text{resp. } H^i(X_{\overline{\eta}}, \Lambda) \to H^i(Y_{\overline{\eta}}, \Lambda))
\]
est bijectif pour $i < m - d - 1$, injectif pour $i = m - d - 1$. Le diagramme commutatif

\begin{tikzcd}
H^i(X_s, \Lambda) \arrow[r, "\varphi^i_X"] \arrow[d] & H^i(X_{\overline{\eta}}, \Lambda) \arrow[d] \\
H^i(Y_s, \Lambda) \arrow[r, "\varphi^i_Y", "\sim"'] & H^i(Y_{\overline{\eta}}, \Lambda)
\end{tikzcd}

(où $\varphi_Y$ est bijectif en vertu de SGA 4 XVI …) montre que \struck{$\varphi^i_X$ est bijectif}
\[
\text{(3.2)} \qquad \varphi^i_X : H^i(X_s, \Lambda) \longrightarrow H^i(X_{\overline{\eta}}, \Lambda)
\]
est bijectif pour $i < m - d - 1$, injectif pour $i = m - d - 1$,
\note{dans le nom du paragraphe et dans (3.1.1)–(3.1.2), « prof. ét. » est l'abréviation de l'auteur ; la ligne sous chaque accolade est écrite au-dessus ou au-dessous de la formule, d'une autre encre.}

\page{140}
\note{p. V.13 de l'auteur. Une grande partie de la page (de « Remarque 3.2 » jusqu'en bas) est barrée de longs traits obliques ; elle est transcrite sous cette réserve.}

en d'autres termes que l'on a
\[
\text{(3.1.3)} \qquad \underline{\Phi}^i_\nu = 0 \quad \text{pour } i \leqslant m - d - 1 .
\]
\struck{Remarque 3.2. Supposons par exemple que $X_s$ soit localement une intersection complète, ce qui \uncertain{implique} \ill{} que $X$ soit loc. une intersection complète, ce qui implique la même chose pour $X_{\overline{\eta}}$. Alors (cf. SGA 1 XIV 5.6) on peut prendre $m \geqslant n$, et (3.1.3) prend la forme}
\[
\text{\struck{(3.2.1)}} \quad \text{\struck{$\underline{\Phi}^i_\nu = 0$ pour $i \leqslant n - d - 1$,}}
\]
\struck{ce qui, compte tenu de 2.7, prend la forme}
\[
\text{\struck{(3.2.2)}} \quad \text{\struck{$\underline{\Phi}^i_\nu = 0$ si $i \notin [n - d, n + d + 1]$.}}
\]
\struck{Si $d = 0$, (ou \uncertain{d'ailleurs} \ill{} de \ill{} $Z$ fini sur $X_0$) cela signifie}
\[
\text{\struck{(3.1.6)}} \quad \text{\struck{$\underline{\Phi}^i_\nu = 0$ si $i \neq n, n+1$.}}
\]
\struck{De \uncertain{même} on a une notion du groupe d'inertie $I$, (3.1.3) \uncertain{signifie} que cette action est triviale pour $H^i(X_{\overline{\eta}}, \Lambda_\nu)$ si $i < m - d - 1$, donc compte tenu de 2.7) si $i \notin [m - d - 1, n + d]$ ; donc si $m = n$ (p.ex. $X_s$ loc. une intersection complète) si $i \notin [n - d - 1, n + d]$, et si de plus $d = 0$ si $i \neq n - 1, n$. On compare avec le \uncertain{corollaire} 2.6, plus fort (car il suffit \uncertain{de supposer} $i \neq n$) mais avec une hypothèse \emph{additionnelle} [$X_\eta$ lisse ou lieu de \ill{} complète] \uncertain{de singularités locales} (des intersections \uncertain{complètes} \add{\uncertain{exactement}} $X_s$ (d'intersection complète)).}
\marginal{Mais dans ce cas \uncertain{on peut améliorer} les conclusions, \uncertain{en remplaçant} $i \leqslant n - d - 1$ par $i \leqslant n - d$.}
\note{l'argument se poursuit au-delà de cette page, hors du lot.}

\end{document}
